\chapter{Strukture podataka}

\index{struktura podataka}

\key{Struktura podataka} je način pohrane
podataka u memoriji računala.
Važno je odabrati prikladnu
strukturu podataka za zadatak,
jer svaka struktura podataka ima svoje
prednosti i nedostatke.
Ključno je pitanje: koje su operacije
efikasne u odabranoj strukturi podataka?

Ovo poglavlje predstavlja najvažnije
strukture podataka iz standardne biblioteke jezika C++.
Dobro je koristiti standardnu biblioteku
kad god je to moguće,
jer će nam uštedjeti puno vremena.
Kasnije u knjizi upoznat ćemo naprednije
strukture podataka koje nisu dostupne
u standardnoj biblioteci.

\section{Dinamički nizovi}

\index{dinamički niz}
\index{vektor}

\key{Dinamički niz} je niz čija se
veličina može mijenjati tijekom izvođenja
programa.
Najpopularniji dinamički niz u C++-u je
struktura \texttt{vector},
koja se može koristiti gotovo kao običan niz.

Sljedeći kod stvara prazan vektor i
dodaje mu tri elementa:

\begin{lstlisting}
vector<int> v;
v.push_back(3); // [3]
v.push_back(2); // [3,2]
v.push_back(5); // [3,2,5]
\end{lstlisting}

Nakon toga elementima možemo pristupati kao u običnom nizu:

\begin{lstlisting}
cout << v[0] << "\n"; // 3
cout << v[1] << "\n"; // 2
cout << v[2] << "\n"; // 5
\end{lstlisting}

Funkcija \texttt{size} vraća broj elemenata u vektoru.
Sljedeći kod prolazi kroz
vektor i ispisuje sve njegove elemente:

\begin{lstlisting}
for (int i = 0; i < v.size(); i++) {
    cout << v[i] << "\n";
}
\end{lstlisting}

\begin{samepage}
Kraći način prolaska kroz vektor je sljedeći:

\begin{lstlisting}
for (auto x : v) {
    cout << x << "\n";
}
\end{lstlisting}
\end{samepage}

Funkcija \texttt{back} vraća posljednji element
u vektoru, a
funkcija \texttt{pop\_back} uklanja posljednji element:

\begin{lstlisting}
vector<int> v;
v.push_back(5);
v.push_back(2);
cout << v.back() << "\n"; // 2
v.pop_back();
cout << v.back() << "\n"; // 5
\end{lstlisting}

Sljedeći kod stvara vektor s pet elemenata:

\begin{lstlisting}
vector<int> v = {2,4,2,5,1};
\end{lstlisting}

Drugi način stvaranja vektora je da zadamo broj
elemenata i početnu vrijednost svakog elementa:

\begin{lstlisting}
// size 10, initial value 0
vector<int> v(10);
\end{lstlisting}
\begin{lstlisting}
// size 10, initial value 5
vector<int> v(10, 5);
\end{lstlisting}

Interna implementacija vektora
koristi običan niz.
Ako se veličina vektora poveća i
niz postane premalen,
alocira se novi niz i svi se
elementi premještaju u novi niz.
Međutim, to se ne događa često pa je
prosječna vremenska složenost operacije
\texttt{push\_back} jednaka $O(1)$.

\index{znakovni niz}

Struktura \texttt{string}
također je dinamički niz koji se može koristiti gotovo kao vektor.
Uz to, za znakovne nizove postoji posebna sintaksa
koja nije dostupna u drugim strukturama podataka.
Znakovne nizove možemo spajati pomoću simbola \texttt{+}.
Funkcija $\texttt{substr}(k,x)$ vraća podniz znakova
koji počinje na poziciji $k$ i ima duljinu $x$,
a funkcija $\texttt{find}(\texttt{t})$ pronalazi poziciju
prvog pojavljivanja podniza znakova \texttt{t}.

Sljedeći kod prikazuje neke operacije nad znakovnim nizovima:

\begin{lstlisting}
string a = "hatti";
string b = a+a;
cout << b << "\n"; // hattihatti
b[5] = 'v';
cout << b << "\n"; // hattivatti
string c = b.substr(3,4);
cout << c << "\n"; // tiva
\end{lstlisting}

\section{Strukture skupova}

\index{skup}

\key{Skup} je struktura podataka koja
održava kolekciju elemenata.
Osnovne operacije nad skupovima su umetanje,
traženje i uklanjanje elementa.

Standardna biblioteka jezika C++ sadrži dvije
implementacije skupa:
Struktura \texttt{set} temelji se na uravnoteženom
binarnom stablu i njezine operacije rade u $O(\log n)$ vremenu.
Struktura \texttt{unordered\_set} koristi raspršivanje,
pa njezine operacije rade u prosječno $O(1)$ vremenu.

Odabir implementacije skupa
često je stvar ukusa.
Prednost strukture \texttt{set}
je što ona održava poredak elemenata
i nudi funkcije koje nisu dostupne
u strukturi \texttt{unordered\_set}.
S druge strane, \texttt{unordered\_set}
može biti efikasniji.

Sljedeći kod stvara skup
koji sadrži cijele brojeve
i prikazuje neke od operacija.
Funkcija \texttt{insert} dodaje element u skup,
funkcija \texttt{count} vraća broj pojavljivanja
elementa u skupu,
a funkcija \texttt{erase} uklanja element iz skupa.

\begin{lstlisting}
set<int> s;
s.insert(3);
s.insert(2);
s.insert(5);
cout << s.count(3) << "\n"; // 1
cout << s.count(4) << "\n"; // 0
s.erase(3);
s.insert(4);
cout << s.count(3) << "\n"; // 0
cout << s.count(4) << "\n"; // 1
\end{lstlisting}

Skup možemo uglavnom koristiti kao vektor,
ali elementima nije moguće pristupati
pomoću zapisa \texttt{[]}.
Sljedeći kod stvara skup,
ispisuje broj elemenata u njemu, a zatim
prolazi kroz sve elemente:
\begin{lstlisting}
set<int> s = {2,5,6,8};
cout << s.size() << "\n"; // 4
for (auto x : s) {
    cout << x << "\n";
}
\end{lstlisting}

Važno svojstvo skupova je
da su svi njihovi elementi \emph{različiti}.
Stoga funkcija \texttt{count} uvijek vraća
ili 0 (element nije u skupu)
ili 1 (element je u skupu),
a funkcija \texttt{insert} nikada ne dodaje
element u skup ako se on
već nalazi u njemu.
Sljedeći kod to ilustrira:

\begin{lstlisting}
set<int> s;
s.insert(5);
s.insert(5);
s.insert(5);
cout << s.count(5) << "\n"; // 1
\end{lstlisting}

C++ sadrži i strukture
\texttt{multiset} i \texttt{unordered\_multiset}
koje inače rade poput struktura \texttt{set}
i \texttt{unordered\_set},
ali mogu sadržavati više primjeraka istog elementa.
Na primjer, u sljedećem kodu sva tri primjerka
broja 5 dodaju se u multiskup:

\begin{lstlisting}
multiset<int> s;
s.insert(5);
s.insert(5);
s.insert(5);
cout << s.count(5) << "\n"; // 3
\end{lstlisting}
Funkcija \texttt{erase} uklanja
sve primjerke elementa
iz multiskupa:
\begin{lstlisting}
s.erase(5);
cout << s.count(5) << "\n"; // 0
\end{lstlisting}
Često želimo ukloniti samo jedan primjerak,
što možemo učiniti ovako:
\begin{lstlisting}
s.erase(s.find(5));
cout << s.count(5) << "\n"; // 2
\end{lstlisting}

\section{Strukture mapa}

\index{mapa}

\key{Mapa} je poopćeni niz
koji se sastoji od parova ključ-vrijednost.
Dok su ključevi u običnom nizu uvijek
uzastopni cijeli brojevi $0,1,\ldots,n-1$,
gdje je $n$ veličina niza,
ključevi u mapi mogu biti bilo kojeg tipa podatka i
ne moraju biti uzastopne vrijednosti.

Standardna biblioteka jezika C++ sadrži dvije
implementacije mape koje odgovaraju
implementacijama skupa: struktura
\texttt{map} temelji se na uravnoteženom
binarnom stablu i pristup elementima
traje $O(\log n)$ vremena,
dok struktura
\texttt{unordered\_map} koristi raspršivanje
pa pristup elementima traje prosječno $O(1)$ vremena.

Sljedeći kod stvara mapu
u kojoj su ključevi znakovni nizovi, a vrijednosti cijeli brojevi:

\begin{lstlisting}
map<string,int> m;
m["monkey"] = 4;
m["banana"] = 3;
m["harpsichord"] = 9;
cout << m["banana"] << "\n"; // 3
\end{lstlisting}

Ako zatražimo vrijednost nekog ključa,
a mapa ga ne sadrži,
ključ se automatski dodaje u mapu s
podrazumijevanom vrijednošću.
Na primjer, u sljedećem kodu
u mapu se dodaje ključ ''aybabtu''
s vrijednošću 0.

\begin{lstlisting}
map<string,int> m;
cout << m["aybabtu"] << "\n"; // 0
\end{lstlisting}
Funkcija \texttt{count} provjerava
postoji li ključ u mapi:
\begin{lstlisting}
if (m.count("aybabtu")) {
    // key exists
}
\end{lstlisting}
Sljedeći kod ispisuje sve ključeve i vrijednosti
u mapi:
\begin{lstlisting}
for (auto x : m) {
    cout << x.first << " " << x.second << "\n";
}
\end{lstlisting}

\section{Iteratori i rasponi}

\index{iterator}

Mnoge funkcije standardne biblioteke jezika C++
rade s iteratorima.
\key{Iterator} je varijabla koja pokazuje
na element u strukturi podataka.

Često korišteni iteratori \texttt{begin}
i \texttt{end} definiraju raspon koji sadrži
sve elemente strukture podataka.
Iterator \texttt{begin} pokazuje na
prvi element strukture podataka,
a iterator \texttt{end} pokazuje na
poziciju \emph{iza} posljednjeg elementa.
Situacija izgleda ovako:

\begin{center}
\begin{tabular}{llllllllll}
\{ & 3, & 4, & 6, & 8, & 12, & 13, & 14, & 17 & \} \\
& $\uparrow$ & & & & & & & & $\uparrow$ \\
& \multicolumn{3}{l}{\texttt{s.begin()}} & & & & & & \texttt{s.end()} \\
\end{tabular}
\end{center}

Uočimo asimetriju u iteratorima:
\texttt{s.begin()} pokazuje na element u strukturi podataka,
dok \texttt{s.end()} pokazuje izvan strukture podataka.
Stoga je raspon definiran iteratorima \emph{poluotvoren}.

\subsubsection{Rad s rasponima}

Iteratori se koriste u funkcijama standardne biblioteke jezika C++
kojima se zadaje raspon elemenata u strukturi podataka.
Obično želimo obraditi sve elemente
strukture podataka, pa funkciji predajemo iteratore
\texttt{begin} i \texttt{end}.

Na primjer, sljedeći kod sortira vektor
pomoću funkcije \texttt{sort},
zatim obrće poredak elemenata pomoću funkcije
\texttt{reverse} i na kraju nasumično miješa poredak
elemenata pomoću funkcije \texttt{random\_shuffle}.

\index{sort@\texttt{sort}}
\index{reverse@\texttt{reverse}}
\index{random\_shuffle@\texttt{random\_shuffle}}

\begin{lstlisting}
sort(v.begin(), v.end());
reverse(v.begin(), v.end());
random_shuffle(v.begin(), v.end());
\end{lstlisting}

Te funkcije možemo koristiti i s običnim nizom.
U tom se slučaju funkcijama predaju pokazivači na niz
umjesto iteratora:

\newpage
\begin{lstlisting}
sort(a, a+n);
reverse(a, a+n);
random_shuffle(a, a+n);
\end{lstlisting}

\subsubsection{Iteratori skupa}

Iteratori se često koriste za pristup
elementima skupa.
Sljedeći kod stvara iterator
\texttt{it} koji pokazuje na najmanji element u skupu:
\begin{lstlisting}
set<int>::iterator it = s.begin();
\end{lstlisting}
Kraći način pisanja tog koda je sljedeći:
\begin{lstlisting}
auto it = s.begin();
\end{lstlisting}
Elementu na koji iterator pokazuje
možemo pristupiti pomoću simbola \texttt{*}.
Na primjer, sljedeći kod ispisuje
prvi element skupa:

\begin{lstlisting}
auto it = s.begin();
cout << *it << "\n";
\end{lstlisting}

Iteratore možemo pomicati operatorima
\texttt{++} (naprijed) i \texttt{--} (natrag),
što znači da se iterator pomiče na sljedeći
ili prethodni element skupa.

Sljedeći kod ispisuje sve elemente
u rastućem poretku:
\begin{lstlisting}
for (auto it = s.begin(); it != s.end(); it++) {
    cout << *it << "\n";
}
\end{lstlisting}
Sljedeći kod ispisuje najveći element skupa:
\begin{lstlisting}
auto it = s.end(); it--;
cout << *it << "\n";
\end{lstlisting}

Funkcija $\texttt{find}(x)$ vraća iterator
koji pokazuje na element čija je vrijednost $x$.
Međutim, ako skup ne sadrži $x$,
iterator će biti \texttt{end}.

\begin{lstlisting}
auto it = s.find(x);
if (it == s.end()) {
    // x is not found
}
\end{lstlisting}

Funkcija $\texttt{lower\_bound}(x)$ vraća
iterator na najmanji element skupa
čija je vrijednost \emph{barem} $x$, a
funkcija $\texttt{upper\_bound}(x)$
vraća iterator na najmanji element skupa
čija je vrijednost \emph{veća od} $x$.
U obje funkcije, ako takav element ne postoji,
povratna vrijednost je \texttt{end}.
Te funkcije nisu podržane u
strukturi \texttt{unordered\_set} koja
ne održava poredak elemenata.

\begin{samepage}
Na primjer, sljedeći kod pronalazi element
najbliži broju $x$:

\begin{lstlisting}
auto it = s.lower_bound(x);
if (it == s.begin()) {
    cout << *it << "\n";
} else if (it == s.end()) {
    it--;
    cout << *it << "\n";
} else {
    int a = *it; it--;
    int b = *it;
    if (x-b < a-x) cout << b << "\n";
    else cout << a << "\n";
}
\end{lstlisting}

Kod pretpostavlja da skup nije prazan
i prolazi kroz sve moguće slučajeve
pomoću iteratora \texttt{it}.
Najprije iterator pokazuje na najmanji
element čija je vrijednost barem $x$.
Ako je \texttt{it} jednak \texttt{begin},
odgovarajući element je najbliži broju $x$.
Ako je \texttt{it} jednak \texttt{end},
najveći element skupa je najbliži broju $x$.
Ako ne vrijedi nijedan od prethodnih slučajeva,
element najbliži broju $x$ je ili
element koji odgovara iteratoru \texttt{it} ili prethodni element.
\end{samepage}

\section{Ostale strukture}

\subsubsection{Bitset}

\index{bitset}

\key{Bitset} je niz
u kojem je svaka vrijednost 0 ili 1.
Na primjer, sljedeći kod stvara
bitset koji sadrži 10 elemenata:
\begin{lstlisting}
bitset<10> s;
s[1] = 1;
s[3] = 1;
s[4] = 1;
s[7] = 1;
cout << s[4] << "\n"; // 1
cout << s[5] << "\n"; // 0
\end{lstlisting}

Prednost korištenja bitsetova je što
zahtijevaju manje memorije od običnih nizova,
jer svaki element bitseta koristi
samo jedan bit memorije.
Na primjer,
ako je $n$ bitova pohranjeno u nizu tipa \texttt{int},
iskoristit će se $32n$ bitova memorije,
dok odgovarajući bitset zahtijeva samo $n$ bitova memorije.
Uz to, vrijednostima bitseta
možemo efikasno manipulirati
bitovnim operatorima, što omogućuje
optimiziranje algoritama pomoću bitsetova.

Sljedeći kod prikazuje drugi način stvaranja gornjeg bitseta:
\begin{lstlisting}
bitset<10> s(string("0010011010")); // from right to left
cout << s[4] << "\n"; // 1
cout << s[5] << "\n"; // 0
\end{lstlisting}

Funkcija \texttt{count} vraća broj
jedinica u bitsetu:

\begin{lstlisting}
bitset<10> s(string("0010011010"));
cout << s.count() << "\n"; // 4
\end{lstlisting}

Sljedeći kod prikazuje primjere korištenja bitovnih operacija:
\begin{lstlisting}
bitset<10> a(string("0010110110"));
bitset<10> b(string("1011011000"));
cout << (a&b) << "\n"; // 0010010000
cout << (a|b) << "\n"; // 1011111110
cout << (a^b) << "\n"; // 1001101110
\end{lstlisting}

\subsubsection{Dvostrani red}

\index{dvostrani red}

\key{Dvostrani red} je dinamički niz
čija se veličina može efikasno
mijenjati na oba kraja niza.
Kao i vektor, dvostrani red nudi funkcije
\texttt{push\_back} i \texttt{pop\_back}, ali
uključuje i funkcije
\texttt{push\_front} i \texttt{pop\_front}
koje nisu dostupne u vektoru.

Dvostrani red možemo koristiti ovako:
\begin{lstlisting}
deque<int> d;
d.push_back(5); // [5]
d.push_back(2); // [5,2]
d.push_front(3); // [3,5,2]
d.pop_back(); // [3,5]
d.pop_front(); // [5]
\end{lstlisting}

Interna implementacija dvostranog reda
složenija je od implementacije vektora
pa je zbog toga dvostrani red sporiji od vektora.
Ipak, i dodavanje i uklanjanje
elemenata traje prosječno $O(1)$ vremena na oba kraja.

\subsubsection{Stog}

\index{stog}

\key{Stog}
je struktura podataka koja nudi dvije
operacije u $O(1)$ vremenu:
dodavanje elementa na vrh
i uklanjanje elementa s vrha.
Moguće je pristupiti samo elementu
na vrhu stoga.

Sljedeći kod prikazuje kako se stog može koristiti:
\begin{lstlisting}
stack<int> s;
s.push(3);
s.push(2);
s.push(5);
cout << s.top(); // 5
s.pop();
cout << s.top(); // 2
\end{lstlisting}
\subsubsection{Red}

\index{red}

\key{Red} također
nudi dvije operacije u $O(1)$ vremenu:
dodavanje elementa na kraj reda
i uklanjanje prvog elementa u redu.
Moguće je pristupiti samo prvom
i posljednjem elementu reda.

Sljedeći kod prikazuje kako se red može koristiti:
\begin{lstlisting}
queue<int> q;
q.push(3);
q.push(2);
q.push(5);
cout << q.front(); // 3
q.pop();
cout << q.front(); // 2
\end{lstlisting}

\subsubsection{Prioritetni red}

\index{prioritetni red}
\index{hrpa}

\key{Prioritetni red}
održava skup elemenata.
Podržane operacije su umetanje te,
ovisno o tipu reda,
dohvaćanje i uklanjanje
najmanjeg ili najvećeg elementa.
Umetanje i uklanjanje traju $O(\log n)$ vremena,
a dohvaćanje traje $O(1)$ vremena.

Iako uređeni skup efikasno podržava
sve operacije prioritetnog reda,
prednost korištenja prioritetnog reda je
u tome što ima manje konstantne faktore.
Prioritetni red obično se implementira pomoću
strukture hrpe koja je mnogo jednostavnija od
uravnoteženog binarnog stabla koje se koristi u uređenom skupu.

\begin{samepage}
Podrazumijevano su elementi u prioritetnom
redu u C++-u sortirani u padajućem poretku,
pa je moguće pronaći i ukloniti
najveći element u redu.
Sljedeći kod to ilustrira:

\begin{lstlisting}
priority_queue<int> q;
q.push(3);
q.push(5);
q.push(7);
q.push(2);
cout << q.top() << "\n"; // 7
q.pop();
cout << q.top() << "\n"; // 5
q.pop();
q.push(6);
cout << q.top() << "\n"; // 6
q.pop();
\end{lstlisting}
\end{samepage}

Ako želimo stvoriti prioritetni red
koji podržava pronalaženje i uklanjanje
najmanjeg elementa,
možemo to učiniti ovako:

\begin{lstlisting}
priority_queue<int,vector<int>,greater<int>> q;
\end{lstlisting}

\subsubsection{Strukture podataka temeljene na politikama}

Prevoditelj (compiler) \texttt{g++} podržava i
neke strukture podataka koje nisu dio
standardne biblioteke jezika C++.
Takve se strukture nazivaju strukturama podataka
\emph{temeljenima na politikama}.
Da bismo koristili te strukture, u kod
moramo dodati sljedeće retke:
\begin{lstlisting}
#include <ext/pb_ds/assoc_container.hpp>
using namespace __gnu_pbds; 
\end{lstlisting}
Nakon toga možemo definirati strukturu podataka \texttt{indexed\_set} koja
je poput strukture \texttt{set}, ali se može indeksirati kao niz.
Definicija za vrijednosti tipa \texttt{int} je sljedeća:
\begin{lstlisting}
typedef tree<int,null_type,less<int>,rb_tree_tag,
             tree_order_statistics_node_update> indexed_set; 
\end{lstlisting}
Sada možemo stvoriti skup ovako:
\begin{lstlisting}
indexed_set s;
s.insert(2);
s.insert(3);
s.insert(7);
s.insert(9);
\end{lstlisting}
Posebnost ovog skupa je što imamo pristup
indeksima koje bi elementi imali u sortiranom nizu.
Funkcija $\texttt{find\_by\_order}$ vraća
iterator na element na zadanoj poziciji:
\begin{lstlisting}
auto x = s.find_by_order(2);
cout << *x << "\n"; // 7
\end{lstlisting}
A funkcija $\texttt{order\_of\_key}$
vraća poziciju zadanog elementa:
\begin{lstlisting}
cout << s.order_of_key(7) << "\n"; // 2
\end{lstlisting}
Ako se element ne pojavljuje u skupu,
dobivamo poziciju koju bi element imao
u skupu:
\begin{lstlisting}
cout << s.order_of_key(6) << "\n"; // 2
cout << s.order_of_key(8) << "\n"; // 3
\end{lstlisting}
Obje funkcije rade u logaritamskom vremenu.

\section{Usporedba sa sortiranjem}

Zadatak je često moguće riješiti
pomoću struktura podataka ili pomoću sortiranja.
Ponekad postoje značajne razlike
u stvarnoj efikasnosti tih pristupa,
koje mogu biti skrivene u njihovim vremenskim složenostima.

Razmotrimo zadatak u kojem su
zadane dvije liste $A$ i $B$
koje obje sadrže $n$ elemenata.
Naš je zadatak izračunati broj elemenata
koji pripadaju objema listama.
Na primjer, za liste
\[A = [5,2,8,9] \hspace{10px} \textrm{i} \hspace{10px} B = [3,2,9,5],\]
odgovor je 3 jer brojevi 2, 5
i 9 pripadaju objema listama.

Izravno rješenje zadatka je
proći kroz sve parove elemenata u $O(n^2)$ vremenu,
no u nastavku ćemo se usredotočiti na
efikasnije algoritme.

\subsubsection{Algoritam 1}

Konstruiramo skup elemenata koji se pojavljuju u $A$,
a nakon toga prolazimo kroz elemente
liste $B$ i za svaki element provjeravamo pripada li
također i listi $A$.
To je efikasno jer se elementi liste $A$
nalaze u skupu.
Uz korištenje strukture \texttt{set},
vremenska složenost algoritma je $O(n \log n)$.

\subsubsection{Algoritam 2}

Nije nužno održavati uređeni skup,
pa umjesto strukture \texttt{set}
možemo koristiti i strukturu \texttt{unordered\_set}.
To je jednostavan način da algoritam učinimo
efikasnijim, jer moramo promijeniti samo
strukturu podataka u pozadini.
Vremenska složenost novog algoritma je $O(n)$.

\subsubsection{Algoritam 3}

Umjesto struktura podataka možemo koristiti sortiranje.
Najprije sortiramo obje liste $A$ i $B$.
Nakon toga prolazimo kroz obje liste
istodobno i pronalazimo zajedničke elemente.
Vremenska složenost sortiranja je $O(n \log n)$,
a ostatak algoritma radi u $O(n)$ vremenu,
pa je ukupna vremenska složenost $O(n \log n)$.

\subsubsection{Usporedba efikasnosti}

Sljedeća tablica pokazuje koliko su efikasni
gornji algoritmi kada $n$ varira i
kada su elementi lista nasumični
cijeli brojevi između $1 \ldots 10^9$:

\begin{center}
\begin{tabular}{rrrr}
$n$ & Algoritam 1 & Algoritam 2 & Algoritam 3 \\
\hline
$10^6$ & $1.5$ s & $0.3$ s & $0.2$ s \\
$2 \cdot 10^6$ & $3.7$ s & $0.8$ s & $0.3$ s \\
$3 \cdot 10^6$ & $5.7$ s & $1.3$ s & $0.5$ s \\
$4 \cdot 10^6$ & $7.7$ s & $1.7$ s & $0.7$ s \\
$5 \cdot 10^6$ & $10.0$ s & $2.3$ s & $0.9$ s \\
\end{tabular}
\end{center}

Algoritmi 1 i 2 jednaki su, osim što
koriste različite strukture skupova.
U ovom zadatku taj izbor ima važan učinak na
vrijeme izvođenja, jer je Algoritam 2
4–5 puta brži od Algoritma 1.

Ipak, najefikasniji je Algoritam 3
koji koristi sortiranje.
On troši samo polovicu vremena u odnosu na Algoritam 2.
Zanimljivo je da je vremenska složenost i
Algoritma 1 i Algoritma 3 jednaka $O(n \log n)$,
ali je unatoč tome Algoritam 3 deset puta brži.
To se može objasniti činjenicom da je
sortiranje jednostavan postupak i da se obavlja
samo jednom na početku Algoritma 3,
a ostatak algoritma radi u linearnom vremenu.
S druge strane,
Algoritam 1 tijekom cijelog algoritma održava
složeno uravnoteženo binarno stablo.
